\subsection{Постановка задачи и результат работы}
По признакам сигналов смартфона требуется определить вид активности человека. Объектом является временное окно измерений --- короткий фрагмент записи. Оно представлено вектором \(x_i\in\mathbb R^d\): списком из \(d\) числовых характеристик. Модель возвращает один класс:
\begin{equation}
\widehat y_i=f_\theta(x_i),\qquad y_i\in\{0,1,2,3,4,5\}.
\end{equation}
Числа обозначают категории: разность между метками не измеряет сходство движений. Устройство модели и её гиперпараметры студент выбирает и объясняет; в качестве сравнимого исходного набора ниже предложены три семейства алгоритмов.

\begin{table}[H]\centering\small
\caption{Словарь меток активности}
\begin{tabularx}{\textwidth}{@{}l l Y@{}}
\toprule
Код & Метка в файле & Смысл \\
\midrule
0 & \file{WALKING} & Ходьба \\
1 & \file{WALKING\_UPSTAIRS} & Ходьба вверх по лестнице \\
2 & \file{WALKING\_DOWNSTAIRS} & Ходьба вниз по лестнице \\
3 & \file{SITTING} & Сидение \\
4 & \file{STANDING} & Стояние \\
5 & \file{LAYING} & Положение лёжа \\
\bottomrule
\end{tabularx}\end{table}

\textbf{Результат сдачи:} таблица сравнения моделей на валидации, обоснованный выбор, итоговый \file{classification\_report} на тесте, матрица ошибок, визуализация сравнения и объяснение влияния гиперпараметров. По завершении работы вы должны уметь различать precision и recall, читать показатели по отдельным классам и объяснять, почему случайное разделение строк бывает недостаточным.

\subsection{От сигналов к признакам}
Описание исходного задания соответствует набору UCI Human Activity Recognition Using Smartphones. В первоисточнике 30 участников выполняли шесть действий со смартфоном на поясе; акселерометр и гироскоп регистрировали сигналы с частотой 50 Гц. Окна длиной 2,56 с содержали 128 отсчётов и перекрывались на 50\%. Из окна извлекались 561 признак во временной и частотной областях~\cite{har}. Число столбцов в учебной CSV-выгрузке проверяется отдельно.

Акселерометр измеряет ускорение, гироскоп --- угловую скорость. Фильтрация ослабляет шум; отделение низкочастотной составляющей помогает различать вклад гравитации и движения тела. Производные сигналов описывают скорость изменения, а евклидова норма объединяет три оси в скалярную величину. Преобразование Фурье показывает, как энергия сигнала распределяется по частотам. В этой работе не требуется заново реализовывать обработку сырых сигналов: используются готовые признаки.

\begin{table}[H]\centering\small
\caption{Как читать обозначения признаков исходного задания}
\begin{tabularx}{\textwidth}{@{}>{\raggedright\arraybackslash}p{0.42\textwidth}Y@{}}
\toprule
Обозначение & Что оно помогает описать \\
\midrule
\file{tBodyAcc}, \file{tGravityAcc}, \file{tBodyGyro} & Ускорение тела, гравитационная составляющая, угловая скорость; \file{t} обозначает временную область. \\
\file{Jerk}, \file{Mag}, \file{XYZ} & Изменение сигнала, его модуль, отдельные оси. \\
Префикс \file{f} & Представление в частотной области после БПФ. \\
\file{mean}, \file{std}, \file{mad}, \file{iqr} & Характеристики положения и разброса; точное определение проверяется по описанию признаков. \\
\file{max}, \file{min}, \file{sma}, \file{energy}, \file{entropy} & Экстремумы, совокупная амплитуда, энергия и характеристика распределения сигнала. \\
\file{arCoeff}, \file{correlation}, \file{angle} & Динамическая структура, связь сигналов, геометрическое соотношение векторов. \\
\file{maxInds}, \file{meanFreq}, \file{bandsEnergy}, \file{skewness}, \file{kurtosis} & Положение частотных компонентов, энергия полос и форма распределения. \\
\bottomrule
\end{tabularx}\end{table}

Не приписывайте признаку смысл только по имени. Например, величина \file{energy} зависит от принятого нормирования. Для последовательности \(a_1,\ldots,a_N\) удобный пример нормированной энергии --- \(N^{-1}\sum_i a_i^2\); это не то же самое, что сумма без деления на \(N\). В отчёте разберите два конкретных столбца своей таблицы.

\subsection{Разделение данных по участникам}
Если соседние окна одного человека попадают и в обучение, и в проверку, модель может воспользоваться индивидуальными особенностями и близостью записей. Результат тогда не обязательно переносится на нового участника. Для данного учебного протокола единица разделения --- \textbf{человек}, а единица классификации --- \textbf{окно сигнала}.

Сохраняйте внешний \file{test.csv} для итоговой проверки. Внутри \file{train.csv} выделите валидацию по идентификатору участника: примерно 20\% участников, фиксированный seed 40. \file{GroupShuffleSplit} определяет долю групп, поэтому доля строк может отличаться от 20\%~\cite{groupsplit}. Проверьте наличие всех шести классов в каждой части. Не выбирайте seed по наилучшей метрике.

\begin{keybox}{Три непересекающихся множества}
Участники обучения и валидации не пересекаются; участники внешнего теста отсутствуют в обеих этих частях. Столбцы \file{Activity} и \file{subject} не поступают в классификатор. Если идентификатора участника нет, нельзя утверждать, что проверена работа на новых людях: сначала уточните состав выгрузки и протокол.
\end{keybox}

Числовые признаки приводят к сопоставимому масштабу, например \(\widetilde x_j=(x_j-\mu_j)/s_j\). Среднее и масштаб вычисляются только на обучающей части. \file{Pipeline} связывает преобразование с моделью и предотвращает случайное применение нового scaler к проверочным объектам~\cite{pitfalls,scaler}.

\subsection{Три семейства классификаторов}
Каждый алгоритм строит своё правило ответа. Сначала разберите его идею словами, затем найдите в коде настройку, которая меняет это правило. Реализовывать алгоритмы библиотек самостоятельно не требуется.
\subsubsection*{Метод ближайших соседей}
Для нового вектора находятся \(k\) ближайших обучающих объектов; класс определяется голосованием. При евклидовом расстоянии
\begin{equation}
\rho(x,z)=\sqrt{\sum_{j=1}^d(x_j-z_j)^2}
\end{equation}
крупный масштаб одного признака может подавить вклад остальных. Поэтому в исходном сравнении KNN используется после стандартизации. Малое \(k\) делает ответ чувствительным к отдельным объектам; большое сильнее сглаживает границу. Существенны также способ взвешивания соседей и метрика расстояния~\cite{knn}.

\subsubsection*{Линейный метод опорных векторов}
При двух признаках линейную границу можно представить прямой на плоскости: точки по разные стороны получают разные ответы. Для большего числа признаков действует та же идея разделения. Метод опорных векторов (SVM) стремится получить границу с широким зазором между классами, допуская штрафуемые нарушения. Параметр \(C\) управляет компромиссом между регуляризацией и штрафом за нарушения разделения: уменьшение \(C\) усиливает регуляризацию --- ограничение сложности подгонки к обучению. Это не универсальное правило улучшения качества. В исходном сравнении используется \file{LinearSVC}; реализация линейного метода опорных векторов~\cite{linearsvc}.

\subsubsection*{Случайный лес}
Дерево последовательно задаёт вопросы о признаках, например «значение больше порога?», и в конце выдаёт класс. Случайный лес (Random Forest) объединяет ответы многих деревьев, обученных со случайным выбором объектов и доступных признаков. Ограничение глубины меняет сложность отдельных деревьев; число деревьев влияет на устойчивость ансамбля и вычислительные затраты. Масштабирование признаков для пороговых разбиений дерева обычно не требуется; для леса в предложенном протоколе scaler не добавляется~\cite{forest}.

\textbf{Гиперпараметры} задаются до обучения конкретной модели: \(k\), \(C\), глубина и число деревьев. Параметры, например коэффициенты линейной модели, определяются при \file{fit}. Изменение гиперпараметров требует нового обучения. Для контроля добавляется \file{DummyClassifier}, всегда предсказывающий самый частый класс обучения. Он показывает результат без анализа признаков: сравнение с ним помогает понять, чему научились остальные модели.

\subsection{Метрики и матрица ошибок}
Пусть \(C_{ij}\) --- число объектов истинного класса \(i\), отнесённых к классу \(j\). Строки матрицы ошибок соответствуют истине, столбцы --- прогнозу. Для класса \(c\):
\begin{align}
TP_c&=C_{cc}, & FP_c&=\sum_{i\ne c}C_{ic}, & FN_c&=\sum_{j\ne c}C_{cj},\\
P_c&=\frac{TP_c}{TP_c+FP_c}, & R_c&=\frac{TP_c}{TP_c+FN_c}, &
F_{1,c}&=\frac{2TP_c}{2TP_c+FP_c+FN_c}.
\end{align}
Precision отвечает на вопрос, какая доля предсказаний класса \(c\) верна. Recall показывает, какая доля настоящих объектов класса \(c\) найдена. F1 объединяет эти характеристики; высокая precision сама по себе не означает высокую полноту.
\begin{equation}
\mathrm{Accuracy}=\frac{\sum_c C_{cc}}{n},\qquad
\mathrm{MacroF1}=\frac{1}{6}\sum_{c=0}^{5}F_{1,c}.
\end{equation}
Macro-F1 придаёт каждому классу одинаковый вес. Weighted-F1 взвешивает классы по их численности; support --- число истинных объектов класса. Для выбора в этой работе используйте macro-F1 валидации, а accuracy и показатели по классам --- для дополнительного анализа~\cite{classification}.

\begin{keybox}{Пример для самопроверки}
Для класса «сидит» модель дала 10 положительных прогнозов: восемь верны, два относятся к другим классам; ещё четыре объекта класса «сидит» пропущены. Тогда precision равна \(8/10=0{,}8\), recall --- \(8/12\approx0{,}667\), F1 --- \(16/22\approx0{,}727\). Это вычислительный пример, не результат работы на данных курса.
\end{keybox}

При нулевом знаменателе метрика не определена обычной формулой. В коде \file{zero\_division=0} задаёт явное соглашение вывода; нулевой support или отсутствие предсказаний следует отметить. Нормирование матрицы по строкам помогает увидеть долю ошибок для каждого истинного класса, а абсолютные числа показывают объём наблюдений.
